\uuid{TNy0}
\titre{Résolution d'équation différentielle du premier ordre}
\auteur{bodin ; croizat ; sacre}
\organisation{exo7}
\licence{CC BY-NC-SA 4.0}
\datecreate{2022-03-08}
\niveau{L1}
\module{Analyse}
\chapitre{Equation différentielle}
\sousChapitre{Résolution d'équation différentielle du premier ordre}
\difficulte{2}
\theme{}
\exo7qcmid{hilisit.equations-differentielles.050}

\contenu{
\begin{question}
\qclassification{225.01}
Soit l'équation différentielle $y'=y+2e^{3x} + 4x e^{3x}$. On recherche une solution particulière sous la forme $f(x) = ax e^{bx}$. Quelles doivent être les valeurs de $a$ et $b$ ?
\begin{answers}
	\bad{$a=4$, $b=3$}
	\good{$a=2$, $b=3$}
	\bad{$a=1$, $b=3$}
	\bad{$a=1$, $b=4$}
\end{answers}
\begin{explanations}
On calcule qu'avec la forme voulue, on a $f'(x) = a e^{bx} + abxe^{bx}$. Ainsi en remplaçant dans l'équation différentielle, on obtient : $a e^{bx} + abxe^{bx} = ax e^{bx} + 2e^{3x} + 4x e^{3x}$, ce qu'on peut écrire $(a+a(b-1)x)e^{bx} = (2+4x)e^{3x}$ pour y voir plus clair. On peut donc identifier dans l'exposant de l'exponentielle que $b=3$. Puis cela donne pour le polynôme qui accompagne les exponentielles $a+2ax = 2 + 4x$, et donc $a=2$.
\end{explanations}
\end{question}
}
